\documentclass[12pt,a4paper,oneside]{article}
\usepackage[brazil]{babel}
\usepackage[utf8]{inputenc}
\usepackage{olymp}


\newlength{\exmpwidouf}
\exmpwidinf=10cm
\exmpwidouf=6cm


\setcounter{problem}{5}

\begin{document}

\begin{problem}{Simplificação de linhas}{linhas}{2}

Em computação gráfica, o processo de simplificação de linhas consiste em reduzir o número de vértices utilizados para representar objetos bidimensionais. Este processo é muito utilizado para preservar visualização em operações de escala e zoom. Ao reduzirmos uma imagem vetorial, linhas e pontos começam a se compactar demasiadamente, prejudicando a clareza da imagem. Para evitar que isso aconteça, pode-se realizar uma simplificação.

Em um dos métodos para simplificação de linhas, calcula-se a área efetiva de cada vértice, dada pela área do triângulo do vértice candidato e seus vizinhos (vértices extremos não possem área efetiva). Na figura abaixo, a região em vermelho representa a área efetiva do vértice $c$:\\[-21pt]

\begin{center}
\includegraphics[scale=0.5]{images/exercise_93_0}
\end{center}
\vspace{-15pt}

Calculadas as áreas efetivas, escolhe-se para remoção o vértice de menor área efetiva.

A área de um triângulo formado por três vértices $A$, $B$ e $C$ pode ser obtida pela seguinte expressão:\\[-18pt]

\begin{equation}
área = \left|\frac{A_x(B_y-C_y) + B_x(C_y-A_y) + C_x(A_y-B_y)}{2}\right|\notag
\end{equation}

Faça um programa que determina o vértice a ser removido de uma linha. Seu programa deverá imprimir o índice do vértice de menor área efetiva, bem como o valor dessa área.\\[-18pt]

%\Notes
%
%Observações, se tiver.

\InputFile

A entrada começa com uma linha contendo um número inteiro $N$, que é o número de vértices da linha. A próxima linha contém $2N$ valores reais, indicando as coordenadas $(X_i,Y_i)$ de cada vértice, na ordem em que aparecem na linha. Ou seja, os dois primeiros valores correspondem às coordenadas $X$ e $Y$ do primeiro vértice, os dois valores seguintes são as coordenadas $X$ e $Y$ do segundo vértice, e assim por diante. Restrições: $3 \leq N \leq 200$.\\[-18pt]

\OutputFile

O programa deverá exibir o índice do vértice que deverá ser removido, isto é, o vértice de menor área efetiva, e o valor da área efetiva deste vértice, com duas casas decimais. Em caso de empate, escolhe entre os de menor área o de menor índice.\\[-18pt]

% \Notes
% Preste atenção aos tipos das variáveis, pois $F(90) \approx 2^{61}$.

\Example

\begin{example}
\exmp{
5
0 2 0.75 1 2 0.5 3 1 4 2
}{
3 0.25
}%
\end{example}

\begin{example}
\exmp{
6
1 2 0 2 1 4 5 6 1 2 5 1
}{
1 1.00
}%
\end{example}


\begin{example}
\exmp{
3
0 0 1 1 2 2
}{
1 0.00
}%
\end{example}


%Comentários sobre o exemplo, se necessário.

\end{problem}

\end{document}
